\section[Diversité des publics]{Vers une bibliothèque numérique orientée usagers : s'adapter à la diversité des publics de Cujas}

La refonte de CujasNum dépasse le cadre d'une simple mise à niveau technique: placer l'usager au cœur de son écosystème constitue un impératif stratégique. La réussite du projet suppose d'abandonner la logique du simple \enquote{répertoire de \ac{PDF}} au profit d'un véritable environnement de travail, capable d'articuler exigences méthodologiques et accessibilité, depuis l'étudiant de master jusqu'au chercheur chevronné \footcite{leblancBibliothequesNumeriquesEnrichies}. L'enjeu est de transformer l'infrastructure statique en un service dynamique soutenant la production scientifique.

L'analyse des entretiens et des retours d'enquête ont mis en évidence différentes typologies de recherche qui imposent des fonctionnalités différenciés: l'enseignant-chercheur expert, le doctorant, l'étudiant et le lecteur éclairé. Ces quatre profils d'usager, utilisant la bibliothèque numérique, ont des méthodes de recherche propres. L'enseignant-chercheur est spécialiste d'une question ou d'une thématique, pratique une recherche de \enquote{fond} sur des sources rares. Ses habitudes de recherche incluent la sérendipité, l'utilisation de l'\ac{OCR} comme d'un outil permettant de consulter le texte d'une façon plus agréable, et télécharge massivement des \ac{PDF} pour alimenter ses propres dossiers \footannexe{\nameref{annexe_A_rapport_audit} pour voir l'entretien réalisé auprès d'Olivier Descamps.}. Le doctorant a des pratiques similaires à celui de l'enseignant-chercheur, mais exige davantage de précision sur la bibliothèque numérique utilisée. Il utilise le numérique pour la préparation pédagogique, nécessitant des outils de capture d'écrans performant pour illustrer ses supports de cours. Pour l'étudiant, la bibliothèque numérique doit être efficace, passant par l'export bibliographique automatisé (Zotero), la rapidité d'accès, des explications précises et des guides d'aide à la prise en main de l'interface et de son contenu. Pour le lecteur éclairé, celui-ci va être attiré par les expositions virtuelles, ou une autre forme de valorisation sur des thématiques d'actualité sur lesquelles il tomberait par hasard pour en apprendre davantage.

Le numérique répond aujourd'hui à une triple nécessité: logistique, scientifique et patrimoniale. Pour les chercheurs, la nouvelle bibliothèque numérique doit leur permettre d'avoir accès au contenu de la bibliothèque mais également aux sources provenant d'autres institutions françaises ou étrangères. Pour les institutions patrimoniales, l'usage du numérique représente un important levier de conservation. Face à la dégradation physique des collections papiers par le temps ou les intempéries, la version numérique offre une pérennité et une \enquote{propreté} que le document physique ne garantit plus. 
 
Actuellement, la bibliothèque numérique ne permet pas une adaptation optimale aux pratiques de la recherche, laissant présumer une compréhension par son public alors que ce n'est pas toujours le cas. Jérôme Dinet l'a expliqué, en déconstruisant le mythe de l'autonomie des étudiants. Face aux développement des nouvelles technologies, il démontre que \enquote{l'aisance technique des nouvelles générations face aux écrans ne garantit pas leur capacité à chercher, filtrer et traiter de l'information académique}\footcite{leblancBibliothequesNumeriquesEnrichies}. Cette fracture d'usage, d'où émergent des experts et des publics non-spécialiste, légitime l'approche d'adapter la bibliothèque numérique en concevant des parcours différenciés et adaptés à chaque niveau de compétence. 
 
Par exemple, actuellement sur CujasNum, les retours des usagers face aux contenus proposés est un mélange d'incompréhension concernant la constitution des collections. En effet, sur le portail de CujasNum, on observe un \enquote{conflit de classification} entre des logiques thématiques et chronologiques. Une autre incompréhension vient des titres portés par les collections, jugés trop précis par rapport au contenu réel de la collection. Cette précision laisse entendre pour le lecteur qu'il trouvera un contenu spécialisé en consultant l'une des collections. Par exemple, la collection intitulée \enquote{Droit romain}, sans médiation éditoriale entraîne une confusion, à l'instar de la présence non-justifiée pour les usagers, d'auteurs du XVIIIe siècle traitants du droit romain \footannexe{\nameref{annexe_A_rapport_audit} pour voir l'entretien réalisé auprès d'une doctorante Floriane Masséna.}. Le retour fait notamment par les doctorants explique qu'en créant une collection \enquote{Droit romain}, on est plutôt en droit d'attendre un mémento des lois, des règles et des principes juridiques régissant la Rome antique. Le manque d'explication derrière les choix faits par les experts métiers entraînent une confusion d'interprétation sur la logique suivie, engendrant une \enquote{peur de rater des sources} \footannexe{\nameref{annexe_A_rapport_audit} pour voir l'entretien réalisé auprès d'une doctorante Floriane Masséna.}. Loin d'être confortable, c'est une des attentes majeures des usagers envers la nouvelle bibliothèque numérique. Il en est une autre qui est de pouvoir retrouver le confort du livre papier qui est recrée numériquement. 